\documentclass[11pt,a4paper]{article}

\usepackage[dutch]{babel}
\usepackage[T1]{fontenc}
\usepackage[utf8]{inputenc}
\usepackage{geometry}
\usepackage{hyperref}
\usepackage{parskip}

\geometry{margin=2.5cm}

\title{Motivatie voor de keuze van de backend-technologie}
\author{ }
\date{ }

\begin{document}

\maketitle

\section{Inleiding}

Voor de ontwikkeling van het digitale archief van de theaterzaal is het belangrijk om een backend-technologie te kiezen die niet alleen vandaag goed werkt, maar ook op lange termijn onderhoudbaar en begrijpelijk blijft.  
Aangezien het systeem in de toekomst mogelijk beheerd zal worden door opvolgers met een beperktere technische achtergrond, is eenvoud en duidelijkheid een cruciale factor.

Op basis van deze vereisten is gekozen voor \textbf{Python in combinatie met het FastAPI-framework} als backend-oplossing.

\section{Waarom Python?}

Python is één van de meest gebruikte programmeertalen ter wereld en staat bekend om zijn leesbaarheid en eenvoud.

De belangrijkste voordelen van Python in dit project zijn:

\begin{itemize}
  \item \textbf{Leesbare syntax}: Python-code lijkt sterk op natuurlijke taal, waardoor de werking van het systeem snel te begrijpen is.
  \item \textbf{Lage instapdrempel}: Nieuwe ontwikkelaars kunnen zich snel inwerken zonder diepgaande voorkennis.
  \item \textbf{Breed ondersteund}: Python wordt actief onderhouden en gebruikt in onderwijs, wetenschap en industrie.
  \item \textbf{Toekomstbestendig}: De kans dat Python in de nabije of verre toekomst verdwijnt is zeer klein.
\end{itemize}


\section{Waarom FastAPI?}

FastAPI is een modern, lichtgewicht webframework voor Python dat speciaal ontworpen is voor het bouwen van web-API’s.

FastAPI sluit goed aan bij de noden van dit project omwille van de volgende redenen:

\subsection{Eenvoud en duidelijkheid}

FastAPI dwingt een duidelijke structuur af:
\begin{itemize}
  \item Elke URL heeft een duidelijke functie
  \item Invoer en uitvoer zijn expliciet gedefinieerd
  \item De code is makkelijk te lezen en te volgen
\end{itemize}

Dit zorgt ervoor dat opvolgers snel begrijpen hoe gegevens worden opgehaald, toegevoegd of aangepast.

\subsection{Automatische documentatie}

Een belangrijk voordeel van FastAPI is dat het automatisch documentatie genereert voor de backend.

Dit betekent dat:
\begin{itemize}
  \item Alle beschikbare functies van het systeem overzichtelijk worden weergegeven
  \item De API getest kan worden via een webinterface
  \item Minder externe documentatie nodig is
\end{itemize}

Dit verlaagt de onderhoudskosten aanzienlijk.

\subsection{Geschikt voor archief-toepassingen}

Het archief vereist vooral:
\begin{itemize}
  \item het ophalen van voorstellingen
  \item het zoeken en filteren van gegevens
  \item het beheren van records
\end{itemize}

FastAPI is zeer geschikt voor deze klassieke \emph{CRUD}-functionaliteit (Create, Read, Update, Delete).


\end{document}
